\chapter{Introduction}

\section{Context and Motivation}
Agents based on large language models increasingly answer questions by querying data through
tools. What such an agent knows about a dataset is whatever the tool tells it: names, types,
descriptions, units, rules. This layer of description decides whether the agent picks the
right table and whether it treats a value correctly.

Two families of description are in use. \emph{Data contracts} describe one dataset at a
time: its schema, its quality guarantees and, in richer variants, its semantics.
\emph{Ontologies} describe a domain: its concepts, the relations between them and the
vocabulary that names them, with datasets attached to that shared model. Contracts are easy
to write and to keep next to the data. Ontologies promise knowledge that crosses datasets
and that can be aligned with standards. \todo{Support both sentences with references.}

This work studies what happens when a system built on data contracts is migrated to an
ontology. The system is TIAGO Dados Abertos, an open-source server that exposes the open
data of the Brazilian power system operator (ONS) through the Model Context Protocol (MCP).
Its 80 datasets cover load, generation, reservoirs, prices and reliability indicators. The
domain is unforgiving: average power is published in MWmed and stored energy in MWmes, and
summing either along time produces a number that looks plausible and is wrong.

The work extends an existing system rather than building a new one. The demand for
evolution came from a practical question raised in the project: whether an ontology would
serve agents better than the contracts already in place. The author is a development intern
at ONS and a member of the team that builds the server and TIAGO PDA, the open-data portal
of the operator; to the best of the team's knowledge, the server is the first MCP server of
the Brazilian power sector.

The server is written in Python 3.12 with the FastMCP framework and answers SQL queries
with DuckDB over Parquet files that ONS publishes on Amazon S3. The study code is also in
Python, the ontology is written in RDF (Turtle), and the agent experiments run on Amazon
Bedrock. \todo{Manual, item 1: one sentence on how the work applies what was learned in
the course (name the disciplines).}

\section{Goal and Research Questions}
The goal of this work, stated with the GQM template \cite{basili1988tame}, is to
\textbf{analyze} the migration of ODCS data contracts to a domain ontology in an MCP server
\textbf{for the purpose of} characterization \textbf{with respect to} migration effort,
mapping fidelity, routing effectiveness, agent effectiveness and context cost
\textbf{from the point of view of} the researcher \textbf{in the context of} the
open-source MCP server TIAGO Dados Abertos, which exposes the open datasets of the
Brazilian power system operator to LLM-based agents. Each aspect gives one research
question, detailed in Section~\ref{sec:goal-rqs}.

\section{Research Method}
The work is a case study following \citeonline{runeson2012case}. The case is the migration
itself, observed in the real server with its real contracts. Evidence comes from four
sources: the contracts and the ontology derived from them, a ledger that records how each
contract construct was represented, two sets of questions written by ONS staff, and the
process records of the migration. Two experiments are embedded in the case: one with the
server's deterministic router and one with an LLM-based agent on Amazon Bedrock.

\section{Contributions}
\begin{itemize}[leftmargin=1.5cm]
\item An ontology for the ONS open-data catalogue whose content is generated from the data
contracts, together with the generator.
\item A construct-level account of the migration: what becomes typed knowledge, what
survives only as text, and what is left behind.
\item An evaluation, on questions written by people, of a router and of an LLM-based agent
fed by contracts, by the ontology, or by both, including context cost.
\item A prototype MCP server that answers from the ontology.
\item Lessons learned from two attempts at the migration.
\end{itemize}

\section{Outline}
Chapter 2 presents the background on LLM-based agents, the Model Context Protocol, data
contracts and ontologies. Chapter 3 discusses related work. Chapter 4 reports the case
study. Chapter 5 concludes with contributions, limitations and future work.
